\documentclass[12pt]{book}

\usepackage{graphicx}
\usepackage{eso-pic}
\usepackage{tcolorbox}
\usepackage{polyglossia}
\setdefaultlanguage[indentfirst=false,forceheadingpunctuation=false]{russian}
\setotherlanguages{english}

\RequirePackage[a4paper, headsep = 0.5 \headsep, left=2.5cm, right=2.1cm, top=2cm, bottom=2.1cm]{geometry}

\setmainfont{URW Bookman}
\newfontfamily\cyrillicfont{URW Bookman}[Script=Cyrillic, Ligatures=TeX]
\setlength{\parindent}{20pt}

\usepackage{enumitem}
\setlist[itemize, 1]{label =\raisebox{-0.3\height}{\scalebox{1.8}{\textbullet}}}

\title{Принимая отчужденность}
\author{Тодд МакГоуэн}
\date{2019}
\begin{document}

\AddToShipoutPictureBG*{\includegraphics[width=\paperwidth]{title.png}}

\null
    \thispagestyle{empty}%
    \addtocounter{page}{-1}%
    \newpage

\tableofcontents

\chapter{Введение}

''Отчужденность'' звучит непривлекательно. Будто что-то разрушительное, чего лучше избегать. Никто не говорит друзьям ''Давайте отчуждаться'', наверняка это не так приятно, как основание сообщества, романтические отношения, общение с приятелями, пицца, игра в теннис или выпивка. Напротив --- мы хотим помочь подростку избежать отчужденности, а отчужденному постороннему --- почувствовать себя уютней. Левые, согласно Карлу Марксу, понимают политическую деятельность как борьбу с отчуждением. Почти всем отчужденность видится проблемой, которую надо решить.

Приговор отчужденности поспешен и забывает о ее освободительном качестве. Исторически, у отчужденности плохая репутация, но ведь именно отчужденность дает возможность выхода за пределы конкретной ситуации. Следовало бы думать о ней гораздо лучше. Увидеть ее достоинства и постараться усилить ее опыт.

Никто открыто не отстаивал отчужденность как проект, но некоторые принимают ее основой свободы. Немецкий идеализм, кульминацией которого стала философия Гегеля, говорит об отчужденности как условии самосознания. Психоанализ воспринимает отчужденность субъекта как непреклонный факт нашего существования и стремится примирить с ним. Экзистенциализм Сартра и Симоны де Бовуар тоже видит в отчужденности субъекта основу свободы. Франц Фанон объединяет все три течения --- Гегеля, психоанализ и экзистенциализм --- и формулирует антиколониальную теорию, выступая против попыток избегать фундаментальной нам всем отчужденности. Впрочем, эти мыслители никогда не используют отчужденность как основное понятие или лозунг, хотя оно объединяет эти разрозненные попытки говорить о свободе. В этой книге отчужденность выдвигается на первый план.

Отчужденность парадоксальна. Она не значит отклонения от исходной идентичности, от того, что мы когда-то имели но потеряли. Отчужденность неизменно присуща нам. Она не отдаленность от первичного, потому что первична. Мы отчуждены в своей субъектности до самой смерти. Смерть --- единственное лекарство от этой болезни\footnote{В момент смерти наша отчужденность проходит --- мы становимся совершенно чуждыми. Смертью субъект совершенно теряется в инаковости, уже не сохраняя ни малейшей частицы себя.}, так что, наверное, не такая уж это и болезнь. Отчужденность --- не то, о чем можно сожалеть или чего лучше избегать, ведь благодаря ней возможно то, что делает существование достойным продолжения --- свобода, равенство, солидарность, способность выходить за пределы ситуации.

Все творение зависит от отчужденности, она выкорчевывает нас из определенности, освобождает от конкретной ситуации. Каждый рождается в специфичном социуме, в котором на него действуют мощные силы, но не сводится к ним. Первичное отчуждение в языке дистанцирует как от других, так и внутренне: к себе теперь можно относиться как к другому. Мы отчуждены от себя и от других, и что бы мы ни делали, нам никогда не преодолеть эту отчужденность. Она основа нашего существования.

Отчужденность --- это отсутствие самоидентичности. Она отдаляет субъекта от условий, в которых он проявляется. Если я не тождествен себе, если я в противоречии с собой, то внешние силы не могут меня полностью определить. Существо, дистанцированное от самого себя, никогда не сможет быть тем, во что его толкают внешние силы. Этот внутренний раскол значит сопротивление детерминированию биологией или обществом. Именно дистанция от самой себя позволяет субъектности действовать против определяющего. Социальное и натуралистическое утрачивают свою власть на субъектом, потому что его отчужденность делает их незначительными.

Я могу не оставаться до конца своих дней жителем маленького американского городка на Среднем Западе, могу бросить вызов определяющим факторам этой ситуации, и все потому что никогда не отождествлен с положением в котором родился или нахожусь. Это дает мне возможность отнестись к ситуации критически, вплоть до отказа от всех идей и ценностей, которые возложило на меня общество. Я могу быть воспитан с энтузиазмом к капиталистическому строю, но стать марксистом. Если меня заставляют быть женственным, я могу стать мужественным. Полюбив Средний Запад Америки, я, возможно, стану работать в Аргентине. Моя неспособность быть именно тем, кто я есть --- моя не-однозначность --- вот что позволяет мне выйти из исходной ситуации. И даже если я остаюсь в ней и принимаю ее, я все равно отношусь к ней как отчужденный. Можно быть чужим и у себя дома. У меня, как и у всех, есть внутренняя чужеродность, мешающая прямому усвоению диктата общества.

Без этой внутренней дистанции человек не был бы субъектом, не смог бы соотноситься с собой. Субъектность порождена и определена этим расколом. Отделенность от самих себя делает нас нами --- неотделимыми от чуждого в нас\footnote{Славой Жижек в \textit{Добавочном Наслаждении} подчеркивает, что отчуждение конститутивно. Мы не целостны изначально, чтобы \textit{затем} отчуждаться. Мы возвращаемся к изначальному через отчуждение. Поэтому ''вне отчужденности никакой жизни нет, нет никакой позитивной основы отчужденности'' (Slavoj Žižek, \textit{Surplus Enjoyment: A Guide for the Non-Perplexed} (London: Bloomsbury, 2022), 39). Верить, что мы были до и вне отчуждения --- ностальгировать по истоку, которого никогда не существовало.}. Чуждое дает нам нашу уникальность.

Эта внутренняя чуждость и есть то, что Зигмунд Фрейд назовет бессознательным\footnote{Для описания своей концепции Фрейд умышленно выберет \textit{unbewußt}, а не распространенный в то время \textit{bewußtlos}. \textit{Bewußtlos} значил бы ''без сознания'', \textit{unbewußt} же значит противоположное сознанию, чуждое его структуре.}. Бессознательное ставит барьер как между субъектом и ситуацией, так и внутри субъекта. Только отчужденные обладают бессознательным. Вернее, чуждый статус бессознательного значит, что это оно владеет нами. Это иное внутри диктует субъекту его действия, направляет усилия и создает убеждения. Бессознательное ориентирует субъектность, потому что выражает ее самораскол\footnote{В психике Фрейд даст бессознательному структурный приоритет. Ускользая от контроля, оно гораздо больше влияет на нас, чем сознательное мышление: ''Бессознательное --- вот истинная психическая реальность'' (Sigmund Freud, \textit{The Interpretation of Dreams}, trans. James Strachey, in The Complete Psychological Works of Sigmund Freud, vol. 5, ed. James Strachey (London: Hogarth Press, 1953 [1900]), 651). Фрейд написал это на заре своей теории, впоследствии он откажется от некоторых своих идей, но никогда не усомнится в приоритете бессознательного.}.

Отчуждение субъектов от себя видно, когда бессознательное заставляет их идти наперекор своим же интересам. Кто-то объедается пончиками и кексами, вместо того, чтобы есть яблоки и брокколи и предотвратить ишемическую болезнь. Другой вместо комфортного, неконфликтного партнера выбирает сложного, болезненного человека, с которым постоянно ссорится. Продолжает искать и употреблять все новые дурманящие вещества, хотя прекрасно понимает, как они опасны и вредны. Хочет профессионального роста, но раз за разом выступает против начальника. Повсюду субъекта поджидают саморазрушительные, ускользающие от сознательного контроля действия.

Каждый сосуществует с идущей вразрез его желаниям силой. Каждый видит, как сам мешает достижению своих целей, как не может контролировать себя. Отсутствие контроля изнурительно, но также основа способности субъекта становиться иным. Верховенство чуждого --- бессознательного --- открывает возможность свободы. Свобода напрямую зависит от подрыва сознательных желаний бессознательными импульсами.

О бессознательном обычно думают как о принуждающем, забирающем свободу. Мы не обдумываем бессознательные действия и не чувствуем, что они выбраны нами. Но свобода --- не плод сознательных размышлений. Это прерывание решимости, определенности. И именно это делает бессознательное. Оно перебивает социальные и психологические определяющих факторы, открывает разрыв в необусловленную ими альтернативную возможность. Избегая сознательного контроля, бессознательное становится точкой свободы субъекта. Мы свободны благодаря своей раздвоенности.

Самотождественности лишена не только субъектность. Ничто не самотождественно, иначе не смогло бы меняться во времени и перемещаться в пространстве\footnote{Жан-Поль Сартр утверждает, что только субъекты отчуждены, а остальные существа тождественны себе. Он радикально различает субъекта, который для-себя, и остальное бытие, которое в-себе: ''Бытие есть. Бытие --- в-себе. Бытие есть то, что оно есть.'' (Jean-Paul Sartre, \textit{Being and Nothingness}, trans. Hazel E. Barnes (New York: Washington Square Press, 1956 [1943]), 29). Сартр здесь не задумывается, как субъект мог стать отчужденным. Будь бытие лишь тем, что оно есть, субъект не возник бы, не проявился из такого чистого царства самотождественности.}. Это касается и животных и растений, и гор и рек. Будь сущность тем, что она есть, она бы никогда не изменилась. Субъект же превращает свою нетождественность в определяющую его структуру, такой процесс и есть отчуждение.

Первое свидетельство отчуждения --- имя. Ни один ребенок не дает его себе сам\footnote{Мама говорит, что я сам выбрал имя: она положила в шляпу несколько листочков бумаги с именами, и я схватил листочек с именем ''Тодд''. Такой метод не избегает навязывания, ведь все равно контролируется взрослым: он выбирает возможные имена, вырезает и кладет в шляпу листочки с ними. Заключенное в присвоении имени отчуждение неизбежно. Затем мама призналась, что ''Тодд'' было написано на каждом из листочков. Впрочем, признание могло быть следствием тогда уже развившейся у нее болезни Альцгеймера, но так или иначе, оно показывает определяющую роль внешнего в присвоении имени.}. Происходит акт насилия, родители или другие взрослые навязывают ребенку произвольное имя. И даже если он потом воспринимает его своим, оно остается чуждым, сошедшим на него извне. Родители могут называть мальчика Эрнестом и надеется, что он вырастет честным и прямодушным. Могут называть девочку Грейс и думать, что она будет приличной и изящной. Иногда используют имена художников, таких как Лэнгстон или Дэшил, пытаясь ассоциацией приобщить детей к их гениальности. Дают имя, связанное с семейной историей, чтобы оно напоминало ребенку о его происхождении, определяло его поведение. Но и ничего не значащее имя остается чуждым названному означающим. Он не выбирал его, и воспринимает его своим только через отчуждение.

Имя это наша дистанция до себя. В каждом языке есть способ идентифицировать себя именем: я могу сказать ''Меня зовут Тодд'', или ''Ich heiße Todd'', или ''Me llamo Todd''. Всякое такое высказывание провозглашает тождество, и этим выявляет внутреннюю дистанцию между говорящим и именем. Так оно выдает, что я и Тодд --- не одно и то же\footnote{В логике первого порядка утверждение тождества исключает различие, потому что различие значило бы категориальную ошибку: смешение идентификации и предикации. На этом замечании Бертран Рассел строит свою критику Гегеля. Гегель ошибается --- утверждает он --- когда говорит об идентичности в различии, об этой фундаментальной отчужденности субъекта, ведь, например, ''Скотт --- автор ''Уэверли'''' предполагает скорее идентификацию, чем предикацию. Критику этой позиции Рассела и защиту Гегеля см. в \textit{Emancipation After Hegel: Achieving a Contradictory Revolution} (New York: Columbia University Press, 2019).}. Самотождественное существо не смогло бы так сказать.

Даже если я переименую себя, имя останется чуждым. Казалось бы, я называю себя сам, не полагаюсь на данное другими. Разве этого недостаточно для преодоления отчужденности? Нет. Называя себя иначе, я действую исходя их доступных возможностей: выбираю из используемых другими имен или комбинирую их. Выбор имени всегда предполагает долг перед внешним. Мое новое имя не только принадлежит другим, но и выбрано с учетом их впечатления. Откуда его не возьми, оно несет печать самораскола.

Неизбежность отчужденности дает некоторую уверенность. И вместо того, чтобы страдать от саморазделения до самой смерти, мы можем манипулировать им и даже усиливать его. Оставление тщетных попыток преодолеть отчужденность, ее принятие, трансформирует ее. Из бремени она становится освобождением от определенности. Переосмысление отчужденности --- важнейшая политическая задача, особенно в эпоху вездесущих схем его якобы преодоления.

Об отчужденности плохо думают, потому что воображают скрывающуюся за ней завершенность. Где-то в затерянном прошлом, в Эдеме, или в утопическом будущем, нас ждет самотождественность, полнота, совершенство. Она всегда ускользает и тем соблазняет еще больше. Такое лишенное всякого отчуждения состояние очень привлекает нас, но лишь пока остается воображаемым. Малейшая попытка его очертить, конкретизировать, воплотить, выдает его полную несостоятельность\footnote{Никто не считает \textit{Возвращенный рай} лучшей поэмой, чем \textit{Потерянный рай}. Отчужденность приятна, потому что ее неприятно пытаться преодолевать. Мильтон невольно демонстирует это разницей в качестве поэм и --- более того --- изображает неотчужденных Бога, Христа, ангелов в \textit{Потерянном раю} совсем не такими привлекательными, как отчужденного Сатану.}.

Нам следует мыслить отчужденность вне ассоциаций с полнотой или самотождественностью. Мы отчуждены, но не от чего-то могущего дополнить нас. Мы и не теряли ничего такого. Помыслив нехватку вне связи с целостностью, мы поменяем политическую значимость отчужденности. Отчужденность не определяется через не-отчужденность.

Отчужденность --- источник всех человеческих страданий, но оно и источник всякого удовлетворения. Отчужденность предоставляет пространство действию вопреки, без него субъект неотличим от сил, породивших его. Заниматься сексом, есть торт или обсуждать немецкий идеализм --- ничего из этого не возможно без отчужденности, потому что без нее нет субъекта.

Если что и отличает говорящих существ от прочих, так это их осознание отсутствия самоидентичности. Через язык отчужденность осознается субъектами неизменно присущей им. Язык создает дистанцию, и это меняет отношение существа к отчужденности. Субъекты означающего могут осознать и принять себя отчужденными, могут найти удовлетворение в этом отсутствии самотождественности, и только отказ рассматривать отчужденность конститутивной составляющей мешает этому.

Чаще мы думаем о своих экзистенциальных и политических проектах как о преодолевающих отчужденность. За редкими исключениями, мы стремимся дать отчужденности бой, искоренить ее. Психотерапевты воображают, что помогают пациентам обрести самоидентификацию. Революционеры стремятся построить общество, в котором жили бы в гармонии с другими. Экоактивисты работают на формой существования, которая была бы в гармонии с природой\footnote{Впрочем, не всегда. Например, эколог Мюррей Букчин, будучи гегельянцем, говорит, что для экологического проекта необходимо принять отчужденность. См. Murray Bookchin, \textit{The Philosophy of Social Ecology: Essays on Dialectical Naturalism} (Chico: AK Press, 2022 [1990]).}. Разительное разнообразие проектов не затрагивает их сути --- избавления от чуждого. Отчужденность непрестанно видится тем, с чем следует бороться, хотя сама борьба возможна только благодаря ней.

Эта книга предлагает совершенно иной взгляд на отчужденность: усилия в ее преодолении значат не радикальный ответ ситуации, а неспособность увидеть конституирующую роль отчужденности. Нам следует оставить попытки преодолеть ее. Отчужденность должна стать нашей экзистенциальной и политической программой. Отчуждение --- это освобождение.

\chapter{Беспокойный субъект}

\section{Продажная идентичность}

Субъект не то же, что присвоенная ему символическая идентичность. Невозможно понять отчуждение, не видя этого базового различия --- между субъектностью и идентичностью. Идентичность складывается из социальных позиций: работа, семейный статус, религиозные пристрастия, политические предпочтения, национальность и т.д.. Как бы легко я себя в них не видел, я никогда не отождествляюсь с ними полностью. Этот недостаток и есть субъектность. Ни одна из идентичностей не может полностью совпасть с ней, как не подгоняй.

Субъектность возникает в отчуждении, в разрыве между субъектом и принимаемой им символической идентичностью. Субъект не способен стать символической идентичностью, потому что соотнесение с ней значит дистанцию от нее. Всегда есть разрыв между мной как субъектом и мной как личностью. Само усилие, порыв соответствовать идентичности выражает расхождение с ней\footnote{Жан-Поль Сартр показывает это в своем рассказе об официанте, ведущем себя так, будто он официант. Пытаясь полностью воплотить символическую идентичность, он неизбежно терпит неудачу: ''Но посмотрите на этого официанта. Его движения быстрые и точные. Чуть слишком точные и чуть слишком быстрые. Он идет к посетителям чуть слишком быстрым шагом, наклоняется вперед чуть слишком нетерпеливо, голос и глаза его выражают чуть слишком заботливый интерес к заказу клиента'' (Jean-Paul Sartre, Being and Nothingness, trans. Hazel E. Barnes (New York: Washington Square Press, 1956 [1943]), 59). Для Сартра вера в то, что человек и есть символическая идентичность --- самообман, и даже обманщик не может преодолеть разрыв между ними. Слишком высокая точность выдает несовершенство.}. Ведь сводись субъектность к символической идентичности, самоидентификация и не понадобилась.

Субъектность неопределима. Кто мы как субъекты --- вопрос без ответа, хотя мы и пытаемся ответить на него: мужчина, женщина, китаец, итальянец, индус, мусульманин, юрист, дантист. Когда я чувствую, что теряю контроль, плотность существующего субъекта, я напоминаю себе, что я американец, или что я сын, или что я профессор --- напоминаю себе, что я есть. Но никакие усилия не снимают дистанцию между субъектностью и идентичностью\footnote{Чем больше прикладывается усилий соответствовать, тем явственней несоответствие. Усилие значит отсутствие того, на что направлено. Если бы я на самом деле соответствовал своей идентичности профессора, мне не пришлось бы стараться одеваться как профессор, говорить как профессор и ходить как профессор.}. Все попытки соответствовать проваливаются\footnote{Субъект не моряк Попай и не может сказать ''Я таков, каков я есть'' (''I Yam what I yam and that's all what I yam'').}.

Символическая идентичность и маскирует, и маркирует вопрос. Путаница неизбежна, потому что субъектности без присвоенной ей символической идентичности не бывает. Идентичность неизбежна, только она ориентирует нас в социальной ситуации.

И в то же время идентичность --- бегство от субъектности. Мы ищем передышку в борьбе за безопасность, надеемся укрыться в идентичности, но она всегда не то, что мы искали. Отношение к ней определяет наше политическое бытие.

Мы можем противостоять символической идентичности отчужденной субъектностью. Кажется, что это бесполезно, но это политический акт --- первый и самый важный, потому что позволяет действовать за пределами символического.

Погружаясь в символическую идентичность, мы избегаем этого противостояния. Избегаем политических вопросов, на которые, кажется, отвечает идентичность. Символическая идентичность значит избегание проблемы свободы. Этот момент подчеркивает Аленка Зупанчич в \textit{Что есть секс?}: ''(освободительная) политика начинается с ''утраты идентичности'', утраты, в которой нет ничего прискорбного''\footnote{Alenka Zupančič, What IS Sex? (Cambridge: MIT Press, 2017), 36}. Идентичность деполитизирует. Всякая идентичность конформна.

Какой бы трансгрессивной не казалась идентичность, субъект, когда ее принимает --- или пытается принять --- подчиняется силам установившей ее социальной власти. Бессознательное желание субъекта всегда проскакивает мимо требований социального авторитета, и в поисках одобрения субъект моделирует себя идентификацией по требуемому образцу. Бессознательное действует вне авторитета, а символическая идентичность ищет его признания. Субъект принимает символическую идентичность, чтобы ее признал социальный авторитет, поэтому она сразу уже им ограничена. И диссонанс между идентичностями не выводит за эти границы, в основе каждой из них все еще лежит социальное признание. Сколь бы радикальной не воображалась идентичность, через нее в конечном счете ищется чье-то одобрение.

Политическое различие между субъектностью и символической идентичностью разоблачается всякий раз, когда возникает скользкая, сложная социальная проблема. Вопрос про отношение к смертной казни, например, провоцирует положиться на символическую идентичность, дать заготовленный ответ, не обдумав все варианты. Принять точку зрения Католицизма и ответить как католик --- что против смертной казни. Упомянуть, что родители воспитали меня в строгом соответствии с законом Десяти Заповедей, так что смертная казнь --- выражение справедливости. Решить, что я либерал, и что, значит, смертная казнь --- слишком суровое наказание. В каждом варианте некоторая символическая идентичность помогает отставить бремя субъектности. Идентичность предлагает способы решения экзистенциальной проблемы субъектности. В любой непонятной ситуации я идентифицируюсь с определенной символической позицией, и получаю ясный ответ.

Переход от субъектности к идентичности закрывает вопросы, на которые нет готовых ответов. Идентичность ориентирует и показывает, как реагировать, поэтому мы ищем в ней убежище. Но как субъект, я должен ответить на экзистенциальные вопросы вне ее руководства. Соответствие ответа некоторой идентичности не оправдывает его.

Символическая идентичность манит в бегство от субъектности\footnote{Принятая идентичность сама может стать в тягость, если я не привык ей соответствовать.}. Поиск идентичности, которая бы полностью нам соответствовала --- это попытка вылечить отчужденность. Субъект видится сводимым к идентичности, когда она затмевает субъектность. Субъект всегда в разладе с собой, но идентичность --- даже если находится в постоянном движении --- создает мираж целостности. Мы можем оставаться в рамках вопроса о субъектности только признавая всякую идентичность неудачной. Такое признание возможно, потому что субъект никогда не соответствует закрепленной за ним идентичности. Контраст с идентичностью выводит субъекта на первый план.

Речь вносит дистанцию в говорящего. Субъектность и есть эта дистанция, эта неспособность совершенно быть собой. Отношение субъекта к самому себе опосредуется означающим, внешним. Означающее и образует форму символической идентичности, и выявляет раскол, заставляющий субъекта искать ответ на вопрос, чем он является. Всякое означающее диалектично: подчинившись ему, существо обретает способность мыслить отрицание себя. В \textit{Прочти мое желание} Джоан Копжек утверждает: ''обозначение неизбежно порождает сомнение, возможность собственного отрицания; научает думать об аннигиляции, полномасштабном разрушении всей нашей означаемой реальности''\footnote{Joan Copjec, \textit{Read My Desire: Lacan Against the Historicists} (Cambridge: MIT Press, 1994), 54–55.}. Без возможности отрицания --- возможности отринуть не только присвоенное значение, но целый означаемый универсум --- нет субъектности. Для отчужденного все данное уязвимо сомнению и может быть отвергнуто, даже утверждать свое означающее он может только через нехватку самотождественности.

Идентифицируя себя как немец, субъект показывает свою нетождественность с немецкостью. Но любая, даже самая прямая идентификация с немецкостью, только лишний раз демонстрирует ее чуждость субъекту. Если бы немецкость действительно не отличалась от него, акт идентификации был излишним и --- более того --- невозможным. Утверждение повторения свидетельствует о различии. Обычно мы не осознаем, как самоидентификация означающими самодистанцирует. Мы используем чужие слова, чтобы описать себя, никто не изобретает означающие. Акт утверждения совпадения --- признак расхождения. Это печать отчужденности.

Марин Ле Пен, лидер французской партии ''Национальное объединение'' заинтересована соответствовать тому, что значит быть француженкой. Для нее Франция --- католическая страна, и должна быть экономически независимой. Традиции Франции, включая язык, отделяют ее от остальной Европы и от мусульманских стран. Но как бы Ле Пен не заботилась о своей французской идентичности, она отличается от нее. Сама ее артикуляция этой идентичности вносит зазор. Никто не привержен традициям Франции так, как Ле Пен, но даже ей приходится трудиться, чтобы укоренить себя в своей французскости.

Самоописание раскрывает самоудаленность. Если бы субъект был идентичен французскости, никакое самоописание не потребовалось бы. Означающее идентичности указывает, что субъект не тот, кем себя называет, что должен отыгрывать свою идентичность, а не просто быть ею. Идентичность субъекта всегда вне его самого. Повторяющимися актами идентификации он принимает и поддерживает ее. Никто не получает идентичность бесплатно.

Множество научно-фантастических фильмов фокусируются на различии между субъектом и его идентичностью. В \textit{Бегущем по лезвию} (Ридли Скотт, 1981), \textit{Вспомнить все} (Пол Верховен, 1990) и \textit{Темный город} (Алекс Пройас, 1998) мы видим, как эта дистанция проявляется в повествовании\footnote{В \textit{Бегущем по лезвию} Рик Декард (Харрисон Форд) --- охотник за головами --- преследует беглых репликантов, искусственно созданных существ, которым для приручения вживили воспоминания. Но понимая, что его собственная идентичность столь же ложна, как идентичности репликантов, Декард освобождается от нее. Он, может быть, тоже репликант, но даже если не так --- в символической идентичности нет надежности. Фильм заканчивается тем, что Декард вступает в романтические отношения с женщиной, зная, что она репликант.}. Персонажи обнаруживают, что самое сокровенное --- воспоминания --- навязано им внешней [злой] силой. Символическая идентичность оказывается чужда субъектности. Дуг Куэйд (Арнольд Шварцнегер) в ''Вспомнить все'' узнает, что раньше был Карлом Хаузером (также Арнольд Шварцнегер) и создал свою нынешнюю личность, чтобы внедриться в восстание на планете Марс и уничтожить его. Все, что он знает о себе --- о своей личности --- ложь, плод манипуляций Хаузера. Личность оказывается внешней манипуляцией. В конце фильма Куэйд вырывается из-под контроля Хаузера, но никакой личности, кроме имплантированной Хаузером у него не остается. Не остается символической идентичности, на которую можно всецело положиться. В Куэйде проявляется то, что касается всякого субъекта, и потому этот персонаж столь притягателен.

Конечно, никто не внедряет в нас фальшивые воспоминания, но, подобно Куэйду, мы не можем полностью отождествиться со своей идентичностью\footnote{Эти фильмы почти всегда показывают отчужденность субъекта через паранойю: внешнее [зло] навязывает ложную идентичность, и тем самым создает раскол в субъекте. Но раскол возникает гораздо раньше и естественней --- на этапе включения субъекта в сферу символического.}. Подобно нему, мы сталкиваемся со своей субъектностью, идущей вразрез с идентичностью. Разве что Куэйд, в отличие от зрителей, не может просто дистанцироваться от своей идентичности, а должен отвергнуть ее как выдумку. Куэйд показывает отчужденную субъектность без сколь-нибудь надежной символической идентичности. Он парадигмальный субъект, совершенно отчужденный, без возможности уткнуться в символическую идентичность. Если мы не как Дуг Куэйд, то тем хуже для нас\footnote{Куэйд, когда встречает лидера восстания Марса --- мутанта Куато, способного читать мысли --- спрашивает, кто он без истинных воспоминаний. На что Куато, как прилежный ученик Сартра, отвечает: ''Ты то, что ты делаешь. Определяют действия, а не память''. Куэйд удовлетворяется таким ответом, и продолжает играть важную роль в революции.}.

Внешний статус идентичности --- отсутствие непосредственной самоидентичности --- порождает бессознательное, действующее вне контроля субъекта. Бессознательное проистекает из разрыва между субъектом и его идентичностью --- между тем, кто он, и тем, что он. Раскол между субъектностью и символической идентичностью соответствует расколу между бессознательным и сознанием. Не всякая символическая идентичность сознательна, но всякая создана в попытке уйти от бессознательного, от чуждого в субъектности.

\section{Аппетит к разрушению}

Есть разрыв между тем, что имеем в виду и тем, что говорим, между намерением и действием, сознанием и бессознательным. Мы --- отчужденные, потому что сказанное нами выражает скорее наше желание, чем то, что мы имели в виду. Наши действия показывают, кто мы, какими бы ни были наши намерения. Бессознательное не тайно, мы всегда демонстрируем его другим своей речью и действиями. И потому другие знают его лучше, чем мы сами: видят его проявление там, где мы видеть не можем. Сами же мы сталкиваемся со своей субъектностью как с чуждым в своей речи и поступках, как с тем, чего не хотели говорить и делать.

Бессознательное желание --- не просто неосознанное желание. Сознание не может принять его как свое, потому что оно нарушает его логику. Это и есть фундаментальный раскол между сознательными и бессознательными желаниями, и этот раскол не преодолеть никакими размышлениями. Даже те, кто совершенно осознает его, все равно поддаются бессознательному желанию, подрывающему сознательные. Субъект постоянно в противоречии с собой. В психоанализе раскол проявляется в негативной терапевтической реакции.

Безуспешность терапии Фрейд объясняет бессознательным. Пациент хочет поправиться, но бессознательное удовлетворяется расстройством и борется за его поддержание. Согласно Фрейду, ''удовлетворение бессознательного чувства вины, пожалуй, самая большая выгода субъекта от болезни, и проявляется в силах, стремящихся сохранить это состояние''\footnote{Sigmund Freud, “The Economic Problem of Masochism” (1924), trans. Joan Riviere, in \textit{The Standard Edition of the Complete Psychological Works of Sigmund Freud}, vol. 19, ed. James Strachey (London: Hogarth Press, 1961), 166.}. Неудачи психоаналитического подхода Фрейда рассказывают о бессознательном гораздо больше, чем его успехи. Пациент, близкий к преодолению своего расстройства, будто удваивает усилия остаться с ним, хотя сознательно хочет выздороветь. Негативная терапевтическая реакция свидетельствует о первичности и превосходстве бессознательного над сознательным в субъекте. Это превосходство и есть отчужденность.

Сознательно субъекты желают преследовать свои интересы и способствовать своему благу. Ставят цели, достижение которых сделает счастливыми. Я хочу жить в комфортной квартире, работать на интересной работе, смотреть увлекательные телепередачи и есть изысканные десерты и уверен, что буду тогда счастлив. Совершенно точно я не хочу быть несчастным, даже если делаю что-то отвратительное --- ем сырые устрицы или занимаюсь борьбой в грязи.

Блез Паскаль был великим теоретиком структуры сознательных желаний. Он отлично понимал важность желания счастья в сознательной жизни субъекта: ''Все без исключения люди стремятся быть счастливыми. Какие бы средства они не использовали --- все стремятся к этой цели. Это мотив каждого поступка каждого человека, даже того, кто решает повеситься''\footnote{Blaise Pascal, \textit{Pensées}, ed. and trans. Roger Ariew (Indianapolis: Hackett, 2005 [1669]), 181.}. Говоря о самоубийце, Паскаль стремится подчеркнуть совершенную повсеместность стремления к счастью: его поиск неизбежен, мы всегда преследуем свои интересы, это наше фундаментальное свойство.

Если ограничиваться сознательным, то анализ Паскаля вполне убедителен: будучи сознательным существом, субъект ищет счастья. Все, что делает сознательное существо --- корыстно, даже самоуничтожение, потому что структура сознания организована вокруг воли субъекта, а не его желания. Бессознательное желание остается слепым пятном в теории Паскаля, и потому он как мыслитель не котируется в наше время. Субъект сознательно желает быть счастливым, но не может избежать бессознательного желания поддерживать в себе отчужденность, что требует продолжающихся актов саморазрушения. Если бы мы были теми, кто мы есть, Паскаль был бы прав. Но разлад в субъектности приводит к действиям, уводящим от счастья. Субъект отчужден от себя и желает сохранить это любой ценой\footnote{Субъект готов платить любую цену, лишь бы сохранить отчужденность, ведь чем больше цена --- тем надежней отчуждение. Ущерб символической идентичности идет на пользу бессознательному желанию, оно процветает в саморазрушающих действиях. Бессознательное питается неудачами сознательного.}.

Уникальность субъекта не в сознательном. Можно представить --- и в научно-фантастических фильмах это показано --- клонирование сознания. Но соотнесенность субъекта с сознательным повторить не получится. Эта соотнесенность и есть бессознательное, основа единичности субъекта. Оно отдаляет нас от сознательного, и благодаря дистанции мы, субъекты, отдельны. Бессознательное не позволяет нам быть теми, кто мы есть, и потому мы выделяемся из природы и культуры. Бессознательное --- добавочное к символическому для получения субъектности\footnote{Сознание отвергает то, что кажется ему отталкивающим, что не соответствует его ценностям, бессознательное же не отвергает ничего. Об этом пишет Фрейд в своем кратком эссе ''Отрицание'': ''В бессознательном нет места ''нет'''' (Sigmund Freud,
''Negation'' (1925), trans. Joan Riviere, in \textit{The Standard Edition of the Complete Psychological Works of Sigmund Freud}, vol. 19, ed. James Strachey (London: Hogarth Press, 1961), 239). Всеприятие заставляет бессознательное принимать и то, что символическая структура скрывает.}. Отчужденность коренится в бессознательном и --- в отличие от сознательных желаний --- несводима ко всему, что на нас влияет.

Субъекты сознательно хотят быть счастливыми, но подсознательно находят удовлетворение в подрыве этой своей воли. Удовлетворение заключается в предотвращении достижения счастья, счастье подобно миражу, что обнажило бы его достижение. Счастье не может быть сейчас. Достигнутость значит, что достигнуто не то, что желалось. Поэтому счастье существует только как обещание счастья. Бессознательное мешает его достижению, тем самым сохраняя его в единственно возможном виде --- в виде мечты.

Изначального удовлетворяющего объекта нет, поэтому бессознательное удовлетворяется созданием и поддержанием отчужденности. Субъект находит удовлетворяющий объект только через его утрату --- утрата позволяет желать. Желание следует рассматривать не как естественное явление, а как структуру, возникающую из взаимосвязи телесных потребностей и социальных требований. Бессознательное желание существует, подрывая сознательное стремление к счастью, постоянно создавая препятствия на пути исполнения сознательных желаний. Бессознательное удовлетворяется в противодействии сознательному.

К каждому высказыванию субъекта бессознательное добавляет значение, которое он не имел в виду, поэтому субъект, будучи пронизан бессознательным, неспособен на лишь констатацию факта. Его бессознательное желание определяет каждое высказывание, и даже самый, казалось бы, невинный комментарий несет его клеймо, которое и заставило субъекта говорить. Когда я говорю соседу, что зловещие тучи предвещают дождь, я не просто предсказываю погоду. Причина этого высказывания кроется за его абстрактным значением, подрывая его невинность. Возможно, я хотел испортить соседу настроение, зная, что он собирался на пикник, и завидуя ему. Или, может, я хочу укрепить связь между нами, показав нас как товарищей по несчастью. Но скорее всего этот банальный комментарий --- способ пережить травму от встречи с соседом. Говоря бессмысленные банальности, я показываю, что мне нечего противопоставить отчуждающей силе его субъектности. Неважно, какой из этих вариантов верен. Бессознательное желание определяет даже самые безобидные замечания, определяет, вообще, все высказывания субъекта. В речи желание неизбежно искажается.

Искажение проявляется и в интерпретации высказываний других. Когда друг говорит мне, что у меня хорошая квартира, я не принимаю это как лишь констатацию факта. Я задаюсь вопросом, почему он это сказал, и отвечаю согласно причине. Возможно, он собирается занять у меня денег, или считает, что обычно моя квартира выглядит грязной, или даже не любит опрятных и собирается разорвать всякие отношения со мной. Желание если и явлено, то во множестве возможностей, а принять утверждение за чистую монету нельзя. Даже не имеющее сознательного манипулятивного намерения, прямое заявление о чистоте моей квартиры несет багаж бессознательного желания. Чтобы говорить, субъект должен хотеть говорить --- и даже молчание подчинено желанию, ведь так же происходит на фоне речи. Если вместо комплимента мой друг промолчит, я все равно буду думать --- теперь над смыслом его молчания. Будучи отчужденным, субъект не может говорить, не извращая сказанное бессознательным. И даже первые, еще не учитывавшие бессознательное, попытки построения теории субъектности укореняли субъектность в отчужденности.

\section{Почему Кант не может быть собой?}

Происхождение термина \textit{субъект} выявляет его внутреннюю связь с отчужденностью. Говоря о субъекте, мыслители всегда говорят об отчужденном субъекте, даже если прямо не указывают это. И они никогда не говорят о субъекте вне его самораскола. \textit{Субъект} о всегда внутренне дистанцированном от себя, в отличие от \textit{я}, \textit{индивида}, \textit{личности} или \textit{человеческого существа} --- эти термины понятней и несут меньше теоретического жаргона, но не подразумевают отчужденность. Конечно, можно уточнить, что речь о раздвоенном я или о раздвоенной личности, но \textit{субъект} уже предполагает это. Не бывает неотчужденного субъекта, всякий субъект --- субъект отчуждения. Это очевидно в том, как термин был первый раз использован в сегодняшнем смысле философом.

Иммануил Кант не изобрел термин \textit{субъект} (\textit{Subjekt} по-немецки), но в \textit{Критике чистого разума} и последующих работах популяризировал его использование в новом смысле. Канта можно по праву считать первым философом субъектности. Начиная с него, \textit{субъект} значит уже не подчиненного --- а именно в этом смысле термин использовали предшественники Канта, например, Барух Спиноза и Готфрид Лейбниц\footnote{Фридрих Ницше будет считать это грубой ошибкой Канта. Согласно Ницше, субъект --- фикция, не реальная сущность, а философов путает грамматика, и они считают субъекта чем-то действительно существующим (в английском и немецком языках \textit{субъект} и \textit{подлежащее} обозначаются одним словом). В \textit{По ту сторону добра и зла} он задает риторический вопрос: ''Не следовало бы нам относиться к субъекту так же, как относимся к сказуемому и дополнению? Не пора ли философам быть выше веры в грамматику?'' (Friedrich Nietzsche, \textit{Beyond Good and Evil: Prelude to a Philosophy of the Future}, trans. Judith Norman (Cambridge: Cambridge University Press, 2002 [1885]), 35). Но Ницше не поясняет, как его концепция сверхчеловека избегает такой грамматической ловушки. Разве нельзя точно так же свести мысль Ницше к грамматике и так отвергнуть?}. После Канта \textit{субъект} значит сущность, негативно соотнесенную с самой собой, не являющуюся тем, чем является\footnote{Это касается не всех --- Луи Альтюссер, например, критикует концепцию субъекта, говоря, что она способствует иллюзии самоидентичности и самообладания, и основывая на этом свою критику идеологической инерпелляции --- процесса, порождающего, согласно Альтюссеру, субъектность.}.

Кант использует концепцию субъектности для отделения сущности, называемой субъектом, от остального мира. Поэтому для Канта субъектность --- всегда отчужденная субъектность, хотя он, конечно, не пишет это явно и не использует термин \textit{отчужденность}\footnote{Кант популяризировал термин \textit{субъект} в философском смысле, Гегель же сделал то же с термином \textit{отчуждение} (Entfremdung). Идею отчуждения обычно связывают с Марксом, но у Гегеля, в отличие от Маркса, она занимает центральное место.}. Согласно Канту, субъект никогда не чувствует себя в мире как дома, потому что конституирует свой мир посредством акта трансцендентального воображения --- без этого не было бы никакого мира, только скопище бессвязных входных данных. Субъект не существует в своем мире, но конституирует его. Субъектность --- условие возможности опыта. Кант, впрочем, не отрицает внешний мир, и старательно --- особенно во втором издании \textit{Критики чистого разума} --- сепарируется от Джорджа Беркли\footnote{Кант добавляет во второе издание \textit{Критики чистого разума} раздел \textit{Опровержение идеализма} исключительно чтобы читатели не спутали его позицию с позицией Беркли.}, но утверждает, что мир получается связным только благодаря трансцендентальному акту субъекта.

Термин \textit{субъект} возникает в начале \textit{Критики чистого разума}, где Кант описывает, как чувственность структурирует наши интуиции. Чувственность --- одна из трех способностей наряду с рассудком и разумом, посредством которых субъект придает воспринимаемому им миру связность. Использование термина \textit{субъект} указывает, что описываемая сущность должна действовать сама на себя, чтобы получить осмысленный опыт. Для Канта субъект --- отделенное от себя и потому способное на себя действовать, существо.

Будучи источником пространства и времени, субъект (в понимании Канта) сам не может быть в мире ни пространственно, ни темпорально. Мир сам по себе не пространственен и не темпорален, эти формы --- вклад субъекта. Согласно Канту, ''Время --- ... лишь субъективное условие нашего созерцания (которое всегда чувственно, т.е. возможно только благодаря воздействию объектов на нас), и само по себе, вне субъекта --- ничто. И все же оно необходимо объективно в отношении всякого явления --- всякой вещи, могущей предстать перед нами в опыте''\footnote{Immanuel Kant, \textit{Critique of Pure Reason}, trans. Paul Guyer and Allen W. Wood (Cambridge: Cambridge University Press, 1998 [1781]), 164.}. Время как форма возникает от субъекта, но ее нельзя отбросить: время структурирует опыт, и поэтому Кант настаивает на его объективности. Но формы пространства и времени также показывают вневременность и внепространственность самого субъекта. Субъектность порождает темпоральность только потому что сама ей не подвержена.

Кантианский субъект есть свое же разделение. Он осознает себя потому что от себя дистанцирован, не тождественен тому, что он есть. Самораскол субъекта особенно заметен в отношении субъектности к свободе и детерминированности. Кант утверждает --- и пытается доказать --- что субъект не свободен \textit{или} детерминирован --- он свободен \textit{и} детерминирован. И дело не в том, что одни его действия свободны, а другие детерминированы --- каждое действие и свободно, и детерминировано. С теоретической точки зрения субъект детерминирован. Но с практической --- как существо в действии --- свободен. И, согласно Канту, оба подхода верны\footnote{Кристина Корсгаард в своем эпохальном исследовании Канта \textit{На пути к царству целей} пишет: ''Возможны два применения разума: теоретическое и практическое. Мы рассматриваем себя как феномены, когда, применяя его теоретически, описываем и объясняем свое поведение; и как ноумены --- когда практически, решая, что делать. Эти две точки зрения не пересекаются, поскольку их цели --- объяснение и решение --- исключают друг друга'' (Christine M. Korsgaard, \textit{Creating the Kingdom of Ends} (Cambridge: Cambridge University Press, 1996), 204). Такая концепция субъектности --- разделенной на теоретическое и практическое --- концепция субъектности как внутренне отчужденной.}.

Только благодаря фундаментальной отчужденности субъект может быть сразу и свободен и детерминирован. Кантианский субъект столь отделен от самого себя, что его можно одинаково успешно мыслить как теоретически, так и практически --- но не теоретически и практически сразу: теория и практика в философии Канта строго несовместимы. Позднее немецкие идеалисты попытаются объединить эти два полюса, что, согласно Канту, рискует исключить фундаментальную для субъекта отчужденность. Фихте, например, будет искать теоретическую перспективу внутри практической, но субъект у него получится уже не столь разделенным, хотя это и исключит для него некоторые проблемы, свойственные философии Канта.

Само разделение философии Канта на теоретическую --- представленную в \textit{Критике чистого разума} --- и практическую --- наиболее известное изложение которой дано в \textit{Критике практического разума} --- символизирует отчужденность субъекта. Эти два направления не просто сосуществуют в его философии, но и чужды и противоречат друг другу. Для теоретического Канта практический Кант --- выход за пределы мыслимого, но для практического Канта теоретический слеп к проблеме свободы субъекта. Нельзя предпочесть одно другому, отчужденность столь важна для Канта, что сама структура его философии проникается ей. По мнению Канта, отчужденность --- не то, что субъект может преодолеть. Потому что это определение субъектности.

\section{Сломанные животные}

Нельзя понять, как работает политика, не принимая во внимание отчужденность. Мы являемся политическими существами не из-за партийной принадлежности, а из-за предвзятого отношения к общественному порядку. Быть политизированным значит не быть нейтральным. Всякий субъект политизирован из-за своей искаженности означающим, своей фундаментальной отчужденности. Отчужденность конституирует политизированность, и нам не отказаться от нее, не убежать от нее в ситуации, кажущиеся деполитизированными. Отчужденность необходима, ни один субъект не может быть сведен к биологии, составляющей его материальную основу\footnote{Согласно Джорджио Агамбену, нацистские лагеря смерти были призваны лишить заключенных в них всякой политической силы, сделать их лишь биологическими существами. В \textit{Останках Освенцима} он пишет о лагере как о месте биовласти, обгладывающей субъекта до его биологических функций, оставляющей от него лишь животное. Он пишет: ''биовласть действует через производство не жизни или смерти, а изменчивого и практически нескончаемого выживания'' (Giorgio Agamben, \textit{Remnants of Auschwitz: The Witness and the Archive}, trans. Daniel Heller-Roazen (New York: Zone Books, 1999 [1998]), 155). Но субъект не может просто выжить --- как бы жестоко на него не действовали репрессивные структуры власти, он остается отчужденным --- а значит и политическим --- существом.}. И дело не в том, что субъектность никак не связана с человеческим животным --- искажение означающим, порождающее субъект, столь фундаментально меняет животное, что должно приниматься отправной точкой в анализе того, что мы делаем. Несмотря на статус биологических существ, мы --- благодаря отчужденности --- можем действовать неестественно. Неестественное склонности в нас превосходят физиологические императивы.

Представляя субъектов естественными и/или культурными сущностями, мы забываем про их отчужденность. Субъективность, понятая как продукт природных или культурных сил, лишена политического бытия --- несводимости к тому, что ее определяет. Натурализация или культурализация субъектности не лишает субъектов политизированности, но делает ее незаметной. Будучи натурализованы или культурализованы, политические решения рассматриваются уже как только выражение врожденного личного интереса, а не как действия свободного субъекта. Но отчужденный субъект --- не корыстное существо, не реагирует так, как ему выгодней, согласно его планам. Субъектности свойственны бессознательные попытки саботировать собственные интересы, находить и создавать себе трудности.

Субъектность не делает человека хозяином остального мира\footnote{Декарт надеялся, что научные исследования, в отличие от чисто спекулятивной философии, ''сделают нас как бы господами и хозяевами природы'' (René Descartes, \textit{Discourse on the Method of Rightly Conducting One’s Reason and Seeking the Truth in the Sciences}, trans. Robert Stoothoff, in The Philosophical Writings of Descartes, vol. 1 (Cambridge: Cambridge University Press, 1985 [1637]), 142–143). Многие комментаторы указывали на опасность такого отношения к природе, Декарт здесь изменяет собственному пониманию субъекта как отдельного от мира и от самого себя, пониманию, которое развивает на протяжении всего \textit{Рассуждения о методе}. Так, ранее в трактате он говорит о равенстве, основываясь на внутренней отчужденности субъекта от места его происхождения.}. Напротив, она отчуждает его от мира, делает его господство над миром невозможным. Из-за отчужденности отношение субъектности к миру всегда на расстоянии и всегда искажено\footnote{Смотреть сквозь символическое --- еще один способ избежать отчужденности. Его в \textit{Реальных Обманах} рассматривает Дженнифер Фридлендер: ''Как субъекты, мы должны бороться с фундаментальной ложью Символического'' (Jennifer Friedlander, \textit{Real Deceptions: The Contemporary Reinvention of Realism} (Oxford: Oxford University Press, 2017), 123). Попытка выйти из блуждания в символической структуре --- это отказ принимать отчужденность. Необходимость фундаментального обмана значит необходимость отчужденности.}.

Субъект в своей основе --- естественное существо, но совершенно извращенное субъектностью. Оно уже неестественно или даже де-естественно. Образ укорененной в природе субъектности на протяжении веков вводил в заблуждение мыслителей. Аристотель, пытаясь определить политическое бытие человечества, обращается к миру природы и нашему к нему отношению. В начале ''Политики'' он дает самое известное определение роли политики в человеческом существовании: ''Человек по своей природе --- политическое животное''\footnote{Aristotle, Politics, trans. B. Jowett, in \textit{The Complete Works of Aristotle}, vol. 2, ed. Jonathan Barnes (Princeton: Princeton University Press, 1984), 1278b.}. Аристотель ошибочно считает источником политизированности наше положение \textit{в} природе, а не отчужденность от нее. Он натурализует политическое бытие говорящего субъекта, что притупляет противоречия, возникающие в обществе и внутри самой субъектности, противоречия, которые сохранялись в человеческой исключительности Платоном\footnote{Сторонники освобождения животных питают слабость к Аристотелю, ведь это принятие преемственности делает его их потенциальным союзником, в отличие от Платона, постулирующего фундаментальный разрыв между животным и человеком. И хотя Аристотель утверждает четкую иерархию среди животных и людей, он все же наделяет животных душой.}.

Политика, согласно Аристотелю --- не борьба противоположностей, а объединение ради общего блага. Он не видит антагонизма между господином и рабом или между мужчиной и женщиной, его политическая философия допускает беспрепятственное угнетение. Понимание человека как нечта естественного, как продолжения природы, прямо связано с присвоением рабам и женщинам статуса второстепенных\footnote{Платон, в отличие от Аристотеля, видит идеальное общество обществом без рабов и ущемления женщин. Он отвергает демократию, считая философов правящим классом в таком обществе, обществе равенства, которое Аристотель не может даже представить.}.

Преемственность человека с миром природы позволяет Аристотелю натурализовать --- и тем самым оправдать --- неравенство внутри политики. Ведь если возможно --- и естественно дозволено --- неравенство в природе, то возможно и дозволено неравенство в человеческом мире. Юридический закон, согласно Аристотелю, действует как естественный закон, описывая уже существующие отношения. Не требует от субъектов действовать против самих себя.

Аристотель так легко принимает неравенство, потому что не осознает отчужденность субъекта от мира природы. Попытка укоренить субъектность в природном всегда ведет к оправданию социального неравенства. Если мы просто природные существа, то равенство невозможно: Флоренс бегает быстрее, чем Инге, Боб умнее Карла, а Карим выше Кальвина. В естественных отношениях доминирует явное неравенство, но, к счастью, мы никогда не остаемся в лишь естественных отношениях друг с другом, чего Аристотель и не учитывает.

Понятие равенства требует говорить о денатурализованном, отчужденном субъекте. Только если политизированность субъекта обеспечена его разрывом с миром природы --- только тогда возможно равенство субъектов. Политическую борьбу не следует рассматривать как лишь поле различных, конкурирующих друг с другом, за продвижение собственных интересов, животных. Именно благодаря денатурализованности субъект способен на антагонистическую борьбу во имя свободы и равенства. То, как мы поймем субъектность, определяет, как мы станем конституировать политическое поле и политические возможности, так что вопрос этот --- не чисто теоретический.

\end{document}
